\section{Boucksom--Jonsson's global theory on projective space}\label{sec:bj22-blueprint}

This chapter collects the material of Boucksom--Jonsson \cite{BJ22} and of their addendum
\cite{BJadd} that is used by Section~3 of \cite{FHX} (the envelope property, Theorem~3.1, for
$X=\PP^n$) and by Lemma~5.10 of \cite{FHX} (\texttt{lem:compactification}, needed in the extension step of
Proposition~\ref{prop:nonflat-norm}). Everything is specialized to $X=\PP^n$, where it simplifies
substantially:
\begin{itemize}
\item test configurations are not used: the graded ideals $J_m$ of \cite[\S3]{FHX} are defined
directly from the obstacle (Definition~\ref{def:bj-Jm});
\item divisorial points are defined through the Gauss extension and discrete valuation rings
essentially of finite type over $k$ (Definition~\ref{def:pn-div}), which is the form in which the
Frobenius trace estimate \cite[Lemma~3.3]{FHX} is applied;
\item the comparison theorem \cite[Theorem~4.22]{BJ22} is proved from Rees valuations of ideals in
polynomial rings (Theorem~\ref{thm:bj-rees}), the only non-elementary input being the finiteness of
normalization (Theorem~\ref{thm:bj-normalization-finite});
\item countable regularization \cite[Corollary~12.18]{BJ22}, which Boucksom--Jonsson derive from
their energy theory, is obtained from second countability of $(\PP^n)^\an$ when $k$ is countable,
which is the standing assumption of Section~5 of \cite{FHX} (Corollary~\ref{cor:bj-countable-regularization}).
\end{itemize}
Nothing in this chapter depends on \cite[Theorem~6.21]{BJ22}.

\begin{convention}\label{conv:bj-pn}
\uses{def:pn-an, def:pn-gauss, def:pn-psh, def:pn-ep}
Throughout, $k$ is algebraically closed with the trivial absolute value, $n\ge1$,
$R=k[x_0,\dots,x_n]$, $R_d$ is the space of forms of degree $d$, and $\mathfrak m=(x_0,\dots,x_n)$.
Points $v\in(\PP^n)^\an$ are the graded normalized semivaluations of Definition~\ref{def:pn-an};
$\sigma(v)$ is the Gauss extension to $R[t]$, so that $\log|\sum_\lambda s_\lambda t^\lambda|_{\sigma(v)}
=\max_\lambda(\log|s_\lambda|_v-\lambda)$, and for an ideal $I\subseteq R[t]$ we write
$\phi_I(v)=\sup_{g\in I}\log|g|_{\sigma(v)}$ (Definition~\ref{def:pn-gauss}). An element of $R[t]$ is
\emph{bihomogeneous} if it is homogeneous in $x$ and a monomial in $t$; an ideal is bihomogeneous if
it is generated by bihomogeneous elements. For $a\in\Q_{\ge0}$, $\mathcal H^{\mathrm{gf}}_\R(a)$,
$\mathcal H^{\mathrm{gf}}_\Q(a)$ and $\mathcal H_\Q(a)$ denote the Fubini--Study functions for
$\mathcal O(a)$ with real weights, with rational weights, and with rational weights and finite values
(Definition~\ref{def:pn-psh}); $\PSH(\theta)$, $\theta\in\R$, is the set of $\theta$-psh functions,
$\mathrm P_\theta$ the psh envelope, $f^\star$ the usc regularization, and $\operatorname{PL}$ the
$\Q$-vector space of rational PL functions, the differences of functions in $\bigcup_a\mathcal H_\Q(a)$
(Definition~\ref{def:pn-ep}). For $0\le j\le n$ we write $y^{(j)}=(x_i/x_j)_{i\ne j}$,
$B_j=k[y^{(j)}][t]$, and $g^{(j)}=g/x_j^{d}\in B_j$ for $g\in R[t]$ homogeneous of $x$-degree $d$; for
a bihomogeneous ideal $I$, $I^{(j)}\subseteq B_j$ is generated by the $g^{(j)}$, $g\in I$ bihomogeneous.
Valuations are written additively; for a discrete valuation ring $V$, $\operatorname{ord}_V$ is its normalized
valuation and $\operatorname{ord}_V(\mathfrak a)=\min_{g\in\mathfrak a}\operatorname{ord}_V(g)$.
\end{convention}

\subsection{Fubini--Study and psh functions}

\begin{lemma}[Operations on Fubini--Study functions]\label{lem:bj-fs-basic}
\lean{EnvelopeProperty.ProjBerk.IsRatFS.mono_class, EnvelopeProperty.ProjBerk.isRatFS_smul_iff, EnvelopeProperty.ProjBerk.IsRatFS.max_const, EnvelopeProperty.ProjBerk.IsFS.exists_antitone_finiteRatFS}
\leanok
\uses{def:pn-psh}
\source{bj22-trivval:prop:FS, bj22-trivval:lem:FSgenfin}
Let $a\le a'$ be nonnegative rationals.
\begin{enumerate}
\item $\mathcal H^{\mathrm{gf}}_\R(a)\subseteq\mathcal H^{\mathrm{gf}}_\R(a')$, and the same holds for
$\mathcal H^{\mathrm{gf}}_\Q$ and $\mathcal H_\Q$.
\item For $c\in\Q_{>0}$, $c\,\mathcal H^{\mathrm{gf}}_\Q(a)=\mathcal H^{\mathrm{gf}}_\Q(ca)$, and similarly for
$\mathcal H^{\mathrm{gf}}_\R$ and $\mathcal H_\Q$.
\item $\mathcal H_\Q(a)$ contains the rational constants and is stable under finite maxima and under
addition of rational constants; $\mathcal H^{\mathrm{gf}}_\Q(a)+\mathcal H^{\mathrm{gf}}_\Q(b)\subseteq
\mathcal H^{\mathrm{gf}}_\Q(a+b)$.
\item Every $\varphi\in\mathcal H^{\mathrm{gf}}_\R(a)$ is the pointwise limit of a decreasing sequence
in $\mathcal H_\Q(a)$.
\end{enumerate}
\end{lemma}

\begin{proof}
(1) Write $\varphi=m^{-1}\max_j(\log|s_j|+\lambda_j)$ with $s_j\in R_{ma}$. Choose $m'\in m\mathbf{Z}_{>0}$ with
$e\coloneqq m'(a'-a)\in\N$. Since $\max_i\log|x_i|=0$, we have $\max_i\log|x_i^e|=0$, hence
$\varphi=m'^{-1}\max_{j,i}\bigl(\log|s_j^{m'/m}x_i^e|+(m'/m)\lambda_j\bigr)$, with forms of degree
$m'a'$. The same formula preserves rationality of the weights and finiteness. (2) and (3) are the
operations of Definition~\ref{def:pn-psh}; constants are $m^{-1}\max_i(\log|x_i^{ma}|+mc)$. (4) Choose
rationals $\lambda_j^{(N)}\downarrow\lambda_j$ and put
$\varphi_N=\max\{m^{-1}\max_j(\log|s_j|+\lambda_j^{(N)}),-N\}$, which lies in $\mathcal H_\Q(a)$ by (3)
and decreases to $\varphi$.
\end{proof}

\begin{lemma}[Decreasing limits of decreasing limits]\label{lem:bj-psh-diagonal}
\lean{EnvelopeProperty.IsDecrNetLimit.of_antitone_of_finite, EnvelopeProperty.IsDecrNetLimit.of_antitone, EnvelopeProperty.IsDecrNetLimit.of_antitone_of_finite_rat}
\leanok
\source{bj22-trivval:lem:PSH}
Let $K$ be a compact space, $S$ a set, $\rho\colon S\to T$ a map to a metric space, and for each
$s\in S$ let $\mathcal F_s$ be a family of continuous functions $K\to\R$ stable under addition of
positive rational constants. For $t\in T$ let $\widetilde{\mathcal F}_t$ be the set of pointwise limits
$g\colon K\to[-\infty,\infty)$ of decreasing nets $f_i\in\mathcal F_{s_i}$ with $\rho(s_i)\to t$. If
$(g_i)_{i\in I}$ is a decreasing net with $g_i\in\widetilde{\mathcal F}_{t_i}$ and $t_i\to t$, then its
pointwise limit $g$ lies in $\widetilde{\mathcal F}_t$.
\end{lemma}

\begin{proof}
For each $i$ choose a decreasing net $(f_{i,j})_{j\in J_i}$, $f_{i,j}\in\mathcal F_{s_{i,j}}$, with limit
$g_i$ and $t_{i,j}\coloneqq\rho(s_{i,j})\to t_i$. Let $A$ be the set of triples $(i,j,\epsilon)$ with
$i\in I$, $j\in J_i$, $\epsilon\in\Q_{>0}$, preordered by $(i,j,\epsilon)\ge(i',j',\epsilon')$ if $i\ge i'$,
$f_{i,j}+\epsilon\le f_{i',j'}+\epsilon'$ on $K$, and
$d(t_{i,j},t_i)+\epsilon\le d(t_{i',j'},t_{i'})+\epsilon'$. Given $\alpha_\nu=(i_\nu,j_\nu,\epsilon_\nu)$,
$\nu=1,2$, choose $i_3\ge i_1,i_2$ and rational $0<\epsilon_3<\min(\epsilon_1,\epsilon_2)$. Since
$g_{i_3}\le g_{i_\nu}\le f_{i_\nu,j_\nu}$ and the $f_{i_\nu,j_\nu}$ are real valued,
$g_{i_3}+\epsilon_3<h\coloneqq\min_\nu(f_{i_\nu,j_\nu}+\epsilon_\nu)$ pointwise. The open sets
$\{f_{i_3,j}+\epsilon_3<h\}$ increase with $j$ and cover $K$, so one of them is $K$; enlarging $j$
further we also get $d(t_{i_3,j},t_{i_3})+\epsilon_3<\min(\epsilon_1,\epsilon_2)$. Then
$(i_3,j,\epsilon_3)$ dominates $\alpha_1,\alpha_2$, so $A$ is directed. The net
$\alpha=(i,j,\epsilon)\mapsto f_{i,j}+\epsilon\in\mathcal F_{s_{i,j}}$ is decreasing. Given $x\in K$ and
$\delta>0$, choose $i_0$ with $g_{i_0}(x)<g(x)+\delta$ (or $<-1/\delta$ if $g(x)=-\infty$) and
$d(t_i,t)<\delta$ for $i\ge i_0$, then $j_0,\epsilon_0$ with $f_{i_0,j_0}(x)+\epsilon_0<g(x)+\delta$
(resp.\ $<-1/\delta$) and $d(t_{i_0,j_0},t_{i_0})+\epsilon_0<\delta$. Every $\alpha\ge(i_0,j_0,\epsilon_0)$
satisfies $f_{i,j}(x)+\epsilon<g(x)+\delta$ and $d(t_{i,j},t)\le d(t_{i,j},t_i)+d(t_i,t)<2\delta$.
\end{proof}

\begin{remark}\label{rmk:bj-psh-diagonal-finite}
\uses{lem:bj-psh-diagonal}
(Formalized variant: \texttt{IsDecrNetLimit.of\_antitone\_of\_finite} is the argument of BJ22 for real-valued families; \texttt{IsDecrNetLimit.of\_antitone} reduces families stable under $\max(\cdot,c)$ and avoiding $+\infty$ to it by truncation, which covers the FS families used for PSH2.) In \cite[Lemma~4.6]{BJ22} (\texttt{lem:PSH}, arXiv version of 2022-03-08) the functions in $\mathcal F_s$ may take the value $-\infty$. The strict
inequality $g_{i_3}+\epsilon_3<h$ used for directedness then fails at points where
$f_{i_1,j_1}=-\infty$, and the argument breaks down (for instance $f_{i_1,j_1}=\log|s|$ and
$f_{i_3,j}=\max\{\log|s|,-j\}$: no $f_{i_3,j}+\epsilon_3$ lies below $\log|s|+\epsilon_1$). The lemma is
therefore stated and used only for real-valued functions; by Lemma~\ref{lem:bj-fs-basic}(4) every
generically finite Fubini--Study function is a decreasing limit of finite-valued ones, which reduces all
applications to finite-valued Fubini--Study functions (proof of Theorem~\ref{thm:bj-psh2}). Only small
positive rational constants are added in the proof, so rational weights are preserved.
\end{remark}

\begin{theorem}[PSH2]\label{thm:bj-psh2}
\lean{EnvelopeProperty.ProjBerk.isPSH_of_antitone, EnvelopeProperty.ProjBerk.isPSH_of_tendsto_of_antitone, EnvelopeProperty.ProjBerk.isPSH_of_antitone_seq}
\leanok
\uses{lem:bj-psh-diagonal, lem:bj-fs-basic}
\source{bj22-trivval:thm:psh12, fhx-envelope-20261028:prop:BJ-package}
\begin{enumerate}
\item A function $u\not\equiv-\infty$ is $\theta$-psh if and only if it is the pointwise limit of a
decreasing net $u_i\in\mathcal H_\Q(a_i)$ with $a_i\in\Q$, $a_i\to\theta$.
\item If $u_i\in\PSH(\theta_i)$ is a decreasing net, $\theta_i\to\theta$, and $u=\lim_iu_i\not\equiv-\infty$,
then $u\in\PSH(\theta)$.
\end{enumerate}
\end{theorem}

\begin{proof}
(1) Let $\mathcal F_b=\mathcal H_\Q(b)$ for $b\in\Q$ and $\rho(b)=b\in\R$. If $u=\inf_i\varphi_i$ with
$\varphi_i\in\mathcal H^{\mathrm{gf}}_\R(b_i)$ decreasing and $b_i\to\theta$, then each $\varphi_i$ lies in
$\widetilde{\mathcal F}_{b_i}$ by Lemma~\ref{lem:bj-fs-basic}(4), and Lemma~\ref{lem:bj-psh-diagonal} gives
$u\in\widetilde{\mathcal F}_\theta$, which is the decreasing net in (1). Conversely such a limit is usc, as an
infimum of continuous functions, and is $\theta$-psh by definition. (2) By (1) each $u_i\in\widetilde{\mathcal F}_{\theta_i}$; by
Lemma~\ref{lem:bj-psh-diagonal}, $u\in\widetilde{\mathcal F}_\theta$, and (1) applies.
\end{proof}

\begin{theorem}[Basic properties]\label{thm:bj-psh-basic}
\lean{EnvelopeProperty.ProjBerk.IsPSH.add, EnvelopeProperty.ProjBerk.IsPSH.smul, EnvelopeProperty.ProjBerk.IsPSH.max, EnvelopeProperty.ProjBerk.IsPSH.add_const, EnvelopeProperty.ProjBerk.isPSH_of_uniform_limit, EnvelopeProperty.ProjBerk.IsPSH.scaleFun, EnvelopeProperty.ProjBerk.IsPSH.exists_common_net, EnvelopeProperty.ProjBerk.IsPSH.convexComp}
\leanok
\uses{thm:bj-psh2, lem:bj-fs-basic, prop:bj-maxprin}
\source{bj22-trivval:thm:PSH, bj22-trivval:lem:twopsh}
Let $\theta,\theta'\ge0$.
\begin{enumerate}
\item $\PSH(\theta)+\PSH(\theta')\subseteq\PSH(\theta+\theta')$ and $t\PSH(\theta)=\PSH(t\theta)$ for $t>0$.
\item $\PSH(\theta)$ is stable under finite maxima, addition of real constants and limits of uniformly
convergent sequences.
\item If $\theta\le\theta'$ then $\PSH(\theta)\subseteq\PSH(\theta')$.
\end{enumerate}
\end{theorem}

\begin{proof}
Use the nets of Theorem~\ref{thm:bj-psh2}(1). Sums: the product net of sums, with classes adding
(Lemma~\ref{lem:bj-fs-basic}(3)). Rational $t$: Lemma~\ref{lem:bj-fs-basic}(2). Real constants $c$: the
approximants plus rationals $c_i\downarrow c$. Maxima of $u,u'\in\PSH(\theta)$: on the product net choose
rationals $b_{ij}\ge\max(a_i,a'_j)$ with $b_{ij}-\max(a_i,a'_j)\to0$; by Lemma~\ref{lem:bj-fs-basic}(1)
both approximants lie in $\mathcal H_\Q(b_{ij})$, and so does their maximum. (3) Constants are
$(\theta'-\theta)$-psh (Definition~\ref{def:pn-psh}), so $u=u+0\in\PSH(\theta')$ by the first part of (1). Real $t$: after subtracting the constant $u(\vtriv)$
(Proposition~\ref{prop:bj-maxprin}) we may assume $u\le0$; for rationals $t_m\uparrow t$ the functions
$t_mu$ decrease to $tu$, and PSH2 applies. Uniform limits: if $u_m\to u$ uniformly, pass to a subsequence with $\sup|u_m-u|\le2^{-m}$; then
$u_m+3\cdot2^{-m}$ decreases to $u$, which is finite wherever the $u_m$ are, and PSH2 applies.
\end{proof}

\begin{proposition}[Monotonicity and maximum principle]\label{prop:bj-maxprin}
\lean{EnvelopeProperty.ProjBerk.IsPSH.antitone, EnvelopeProperty.ProjBerk.IsPSH.iSup_eq, EnvelopeProperty.ProjBerk.IsPSH.rpow_convex, EnvelopeProperty.ProjBerk.IsPSH.apply_triv_eq_zero}
\leanok
\uses{thm:bj-psh2}
\source{bj22-trivval:prop:pshconvex, bj22-trivval:lem:maxprin}
Order $(\PP^n)^\an$ by $v\le w$ if $|f|_w\le|f|_v$ for all $f\in R$. Every $u\in\PSH(\theta)$ is
decreasing for this order, $\vtriv$ is its least element, and $\sup u=u(\vtriv)\in\R$.
\end{proposition}

\begin{proof}
A function $\log|s|$, $s\in R_d$, is decreasing by definition of the order, hence so are Fubini--Study
functions and their decreasing limits. Since $|f|_v\le1=|f|_{\vtriv}$ for $0\ne f\in R$ homogeneous and $v$
is graded, $\vtriv\le v$ for all $v$. Thus $u\le u(\vtriv)$, which is finite since $u\not\equiv-\infty$
and $u<+\infty$.
\end{proof}

\begin{lemma}[Envelopes]\label{lem:bj-env-basic}
\lean{EnvelopeProperty.ProjBerk.envelope_add_const, EnvelopeProperty.ProjBerk.envelope_mul}
\leanok
\uses{thm:bj-psh-basic}
\source{fhx-envelope-20261028:subsec:prelim-envelopes}
For $\theta\ge0$ and bounded $\phi$: $\mathrm P_\theta$ is increasing,
$\mathrm P_\theta(\phi+c)=\mathrm P_\theta(\phi)+c$ for $c\in\R$, and
$\mathrm P_{t\theta}(t\phi)=t\,\mathrm P_\theta(\phi)$ for $t>0$.
\end{lemma}

\begin{proof}
$u\mapsto u+c$ is a bijection of $\PSH(\theta)$ and $u\mapsto tu$ a bijection
$\PSH(\theta)\to\PSH(t\theta)$ (Theorem~\ref{thm:bj-psh-basic}); both preserve the constraint $u\le\phi$
in the corresponding form.
\end{proof}

\subsection{Regularization}

\begin{theorem}[Regularization]\label{thm:bj-regularization}
\lean{EnvelopeProperty.ProjBerk.IsPSH.exists_finiteRatFS_net, EnvelopeProperty.ProjBerk.IsPSH.exists_finiteRatFS_net_of_rat}
\leanok
\uses{thm:bj-psh2, thm:bj-psh-basic, prop:bj-maxprin, lem:bj-psh-diagonal}
\source{bj22-trivval:thm:pshample}
Let $\theta\ge0$ and $u\in\PSH(\theta)$.
\begin{enumerate}
\item $u$ is the limit of a decreasing net $u_\alpha\in\mathcal H_\Q(a_\alpha)$ with $a_\alpha\in\Q$,
$a_\alpha>\theta$ and $a_\alpha\to\theta$.
\item If $\theta=a\in\Q_{>0}$, then $u$ is the limit of a decreasing net in $\mathcal H_\Q(a)$.
\end{enumerate}
\end{theorem}

\begin{proof}
(1) Let $u_i\in\mathcal H_\Q(b_i)$ be as in Theorem~\ref{thm:bj-psh2}(1). On $I\times\N$ choose rationals
$a_{i,N}\in(\max(b_i,\theta),\max(b_i,\theta)+1/N)$; then $u_i\in\mathcal H_\Q(a_{i,N})$
(Lemma~\ref{lem:bj-fs-basic}(1)), the net $(i,N)\mapsto u_i$ decreases to $u$, and $a_{i,N}\to\theta$.
(2) Choose a rational $c\ge u(\vtriv)=\sup u$ (Proposition~\ref{prop:bj-maxprin}) and put $u'=u-c\le0$.
Take $u'_\alpha$, $a_\alpha$ as in (1) for $u'$. For rational $\tau>1$, eventually $a_\alpha\le\tau a$, so
$\tau^{-1}u'_\alpha\in\mathcal H_\Q(a)$ (Lemma~\ref{lem:bj-fs-basic}(1),(2)), and these decrease to
$\tau^{-1}u'$. As $\tau\downarrow1$, $\tau^{-1}u'$ decreases to $u'$ because $u'\le0$.
Lemma~\ref{lem:bj-psh-diagonal}, with $T$ a point and $\mathcal F=\mathcal H_\Q(a)$, gives a decreasing net
in $\mathcal H_\Q(a)$ with limit $u'$; add the rational constant $c$.
\end{proof}

\begin{lemma}[Countable subnets]\label{lem:bj-lindelof}
\lean{EnvelopeProperty.exists_monotone_iInf_eq}
\leanok
Let $Z$ be a second countable space and $(u_i)_{i\in I}$ a decreasing net of usc functions
$Z\to[-\infty,\infty)$ with limit $u$. There are indices $i_1\le i_2\le\cdots$ such that
$u=\inf_mu_{i_m}$; moreover the $i_m$ may be chosen above any given sequence of indices.
\end{lemma}

\begin{proof}
For $q\in\Q$ the open set $\{u<q\}$ is the union of the open sets $\{u_i<q\}$. Second countable spaces
are hereditarily Lindel\"of, so $\{u<q\}=\bigcup_{i\in I_q}\{u_i<q\}$ for a countable $I_q\subseteq I$.
Enumerate $\bigcup_qI_q=\{i'_1,i'_2,\dots\}$ and choose, by directedness, $i_m\ge i_{m-1},i'_m$ (and
above the given sequence). If $\inf_mu_{i_m}(z)>q>u(z)$ for some $z$ and $q\in\Q$, then
$u_{i'}(z)<q$ for some $i'=i'_m\in I_q$, and $u_{i_m}(z)\le u_{i'_m}(z)<q$, a contradiction.
\end{proof}

\begin{corollary}[Countable regularization]\label{cor:bj-countable-regularization}
\lean{EnvelopeProperty.ProjBerk.IsPSH.exists_finiteRatFS_seq, EnvelopeProperty.ProjBerk.IsPSH.exists_finiteRatFS_seq_of_rat}
\leanok
\uses{thm:bj-regularization, lem:bj-lindelof, lem:polydisc-compact}
\source{bj22-trivval:cor:countreg}
Assume $k$ countable. In Theorem~\ref{thm:bj-regularization} the nets can be replaced by sequences.
\end{corollary}

\begin{proof}
$(\PP^n)^\an$ is a subspace of the product $\R_{\ge0}^{R}$ over the countable set $R$, hence second
countable. Apply Lemma~\ref{lem:bj-lindelof} to the net of Theorem~\ref{thm:bj-regularization}, choosing
the $i_m$ above indices beyond which $|a_\alpha-\theta|<1/m$.
\end{proof}

\begin{remark}\label{rmk:bj-countable-deviation}
\uses{cor:bj-countable-regularization}
\cite[Corollary~12.18]{BJ22} holds for arbitrary $k$, but its proof uses the energy functionals and the
strong topology of \cite[\S\S7--12]{BJ22}. Countability of $k$ is the convention of Section~5 of
\cite{FHX}, where sequences are needed (regularizable functions and locally psh functions are defined by
sequences); Section~3 of \cite{FHX} takes decreasing limits only of sequences of Fubini--Study functions (by the definition of psh
functions, \cite[Definition~4.1]{BJ22}; by PSH2 in revision \texttt{832e086}) and uses no countable regularization.
\end{remark}

\subsection{PL functions and the envelope criterion}

\begin{lemma}[PL functions form a separating lattice]\label{lem:bj-pl-lattice}
\lean{EnvelopeProperty.ProjBerk.IsRatPL.max, EnvelopeProperty.ProjBerk.IsRatPL.min, EnvelopeProperty.ProjBerk.exists_isRatPL_ne}
\leanok
\uses{lem:bj-fs-basic, def:pn-ep}
\source{bj22-trivval:lem:FSdense}
$\operatorname{PL}$ is a $\Q$-vector space of continuous real functions, contains the rational constants,
is stable under finite maxima and minima, and separates the points of $(\PP^n)^\an$.
\end{lemma}

\begin{proof}
For $\varphi_\nu,\psi_\nu$ finite rational Fubini--Study functions,
$\max(\varphi_1-\psi_1,\varphi_2-\psi_2)=\max(\varphi_1+\psi_2,\varphi_2+\psi_1)-(\psi_1+\psi_2)$, and the two
functions in the maximum lie in a common $\mathcal H_\Q(a)$ by Lemma~\ref{lem:bj-fs-basic}(1),(3); minima
follow by negation. If $v\ne w$, there is $f\in R$ with $|f|_v\ne|f|_w$, and since both points are graded
there is a homogeneous such $f=s\in R_d$. Then $\max\{\log|s|,-N\}=\max\{\log|s|,\max_i\log|x_i^d|-N\}$
lies in $\mathcal H_\Q(d)$ and separates $v$ and $w$ for $N$ large.
\end{proof}

\begin{theorem}[Density of PL functions]\label{thm:bj-pl-dense}
\lean{EnvelopeProperty.ProjBerk.exists_isRatPL_abs_sub_lt, EnvelopeProperty.ProjBerk.dense_setOf_isRatPL}
\leanok
\uses{lem:bj-pl-lattice}
\source{bj22-trivval:thm:PLdense}
$\operatorname{PL}$ is dense in $C^0((\PP^n)^\an,\R)$ for uniform convergence.
\end{theorem}

\begin{proof}
The uniform closure $\overline{\operatorname{PL}}$ is a closed sublattice containing the real constants
and stable under real scalar multiplication (limits of rational ones). If $v\ne w$ and $\psi\in\operatorname{PL}$
separates them, then for any $a,b\in\R$ some $\alpha\psi+\beta\in\overline{\operatorname{PL}}$ takes the
values $a,b$ at $v,w$. The lattice version of the Stone--Weierstrass theorem on the compact space
$(\PP^n)^\an$ gives $\overline{\operatorname{PL}}=C^0$.
\end{proof}

\begin{lemma}[Approximation of usc functions from above]\label{lem:bj-usc-pl}
\lean{EnvelopeProperty.ProjBerk.exists_isRatPL_gt, EnvelopeProperty.ProjBerk.exists_antitone_isRatPL_net, EnvelopeProperty.ProjBerk.le_of_forall_isRatPL_gt}
\leanok
\uses{thm:bj-pl-dense}
Let $h\colon(\PP^n)^\an\to[-\infty,\infty)$ be usc. The set $\{\psi\in\operatorname{PL}:\psi>h\}$ is
nonempty, directed downwards, and its pointwise infimum is $h$.
\end{lemma}

\begin{proof}
$h$ is bounded above by compactness, so a large constant lies in the set; it is stable under minima by
Lemma~\ref{lem:bj-pl-lattice}. Let $x$ be a point, $r>h(x)$ real and $\delta>0$. The set $F=\{h\ge r\}$ is
closed and misses $x$; Urysohn's lemma gives a continuous $f\colon(\PP^n)^\an\to[0,1]$ with $f(x)=0$ and
$f=1$ on $F$. With a real $C\ge\max(\sup h,r)$, the continuous function $g=r+(C+1-r)f$ satisfies
$g\ge r>h$ off $F$, $g=C+1>h$ on $F$, and $g(x)=r$. By Theorem~\ref{thm:bj-pl-dense} choose $\psi\in\operatorname{PL}$
with $|\psi-(g+2\delta)|<\delta$; then $\psi>h$ and $\psi(x)<r+3\delta$.
\end{proof}

\begin{lemma}[Approximation of classes]\label{lem:bj-envppty}
\lean{EnvelopeProperty.ProjBerk.HasEnvelopeProperty.of_tendsto, EnvelopeProperty.ProjBerk.hasEnvelopeProperty_of_forall_rat_gt}
\leanok
\uses{thm:bj-psh2, thm:bj-psh-basic}
\source{bj22-trivval:lem:envppty}
If $\theta\ge0$ is the limit of classes $\theta_j\ge\theta$ having the envelope property, then $\theta$ has
the envelope property.
\end{lemma}

\begin{proof}
Let $(u_i)$ be a nonempty family in $\PSH(\theta)$ bounded above, $u=(\sup_iu_i)^\star$. By
Theorem~\ref{thm:bj-psh-basic}(3), $u_i\in\PSH(\theta_j)$, so $u\in\PSH(\theta_j)$ for all $j$. The constant
sequence $u$ is decreasing, $\theta_j\to\theta$, and $u\ge u_{i_0}\not\equiv-\infty$; PSH2 gives
$u\in\PSH(\theta)$.
\end{proof}

\begin{proposition}[Envelope property via PL obstacles]\label{prop:bj-envelope-criterion}
\lean{EnvelopeProperty.ProjBerk.HasEnvelopeProperty.isPSH_envelope, EnvelopeProperty.ProjBerk.hasEnvelopeProperty_of_forall_isRatPL, EnvelopeProperty.ProjBerk.hasEnvelopeProperty_iff_forall_isRatPL}
\leanok
\uses{lem:bj-usc-pl, thm:bj-psh2, lem:bj-env-basic}
\source{fhx-envelope-20261028:prop:ample-envelopes, bj22-addendum:lem:contenv}
For $\theta\ge0$ the following are equivalent: (i) $\theta$ has the envelope property; (ii)
$\mathrm P_\theta(\phi)\in\PSH(\theta)$ for every $\phi\in\operatorname{PL}$.
\end{proposition}

\begin{proof}
This is \cite[Lemma~1.3]{BJadd}. (i)$\Rightarrow$(ii): the family $\{u\in\PSH(\theta):u\le\phi\}$ contains
the constant $\min\phi$ and is bounded above by $\max\phi$, so $\mathrm P_\theta(\phi)^\star$ is $\theta$-psh;
it is $\le\phi$ since $\phi$ is continuous, hence $\mathrm P_\theta(\phi)^\star\le\mathrm P_\theta(\phi)$, and
the two agree. (ii)$\Rightarrow$(i): let $(u_i)$ be bounded above in $\PSH(\theta)$, $h=(\sup_iu_i)^\star$.
For $\psi\in\operatorname{PL}$ with $\psi>h$ we have $u_i\le\mathrm P_\theta(\psi)$, hence
$h\le\mathrm P_\theta(\psi)\le\psi$, as $\mathrm P_\theta(\psi)$ is usc. By Lemma~\ref{lem:bj-usc-pl} the
functions $\mathrm P_\theta(\psi)$, indexed by the downward directed set of such $\psi$, form a decreasing
net of $\theta$-psh functions with infimum $h\not\equiv-\infty$; PSH2 gives $h\in\PSH(\theta)$.
\end{proof}

\subsection{Fubini--Study envelopes and the ideals \texorpdfstring{$J_m$}{Jm}}

\begin{lemma}[Fubini--Study envelopes]\label{lem:bj-fs-envelope}
\lean{EnvelopeProperty.ProjBerk.envelope_eq_iSup_finiteRatFS, EnvelopeProperty.ProjBerk.envelope_coe_eq_iSup_finiteRatFS}
\leanok
\uses{thm:bj-regularization, lem:bj-fs-basic}
\source{bj22-trivval:lem:envlsc, fhx-envelope-20261028:prop:ample-envelopes, bj22-addendum:lem:FSenv}
Let $a\in\Q_{>0}$ and $\phi\colon(\PP^n)^\an\to\R$ bounded and lsc. Then
$\mathrm P_a(\phi)=\sup\{\psi\in\mathcal H_\Q(a):\psi\le\phi\}$.
\end{lemma}

\begin{proof}
This is \cite[Lemma~1.5]{BJadd} for an ample class. The inequality $\ge$ holds because Fubini--Study
functions are psh. Let $\tau\in\PSH(a)$ with $\tau\le\phi$, and $\epsilon\in\Q_{>0}$. Take
$\tau_\alpha\in\mathcal H_\Q(a_\alpha)$ as in Theorem~\ref{thm:bj-regularization}(1). The sets
$\{\tau_\alpha<\phi+\epsilon\}$ are open (as $\phi-\tau_\alpha$ is lsc), increase, and cover the compact
space, so $\tau_\alpha<\phi+\epsilon$ for all large $\alpha$. Put $M=\sup|\phi|+1$ and choose rationals
$\delta_\alpha\ge|a/a_\alpha-1|M$. Then $(a/a_\alpha)\tau_\alpha-\epsilon-\delta_\alpha\in\mathcal H_\Q(a)$
(Lemma~\ref{lem:bj-fs-basic}(2),(3)) and is $\le\phi$. Hence $(a/a_\alpha)\tau_\alpha\le\mathrm Q+\epsilon+\delta_\alpha$
with $\mathrm Q$ the right-hand side; letting $\alpha\to\infty$ at each point and then $\epsilon\to0$ gives
$\tau\le\mathrm Q$.
\end{proof}

\begin{definition}[The ideals $J_m$]\label{def:bj-Jm}
\lean{EnvelopeProperty.ProjBerk.obstacleSpace, EnvelopeProperty.ProjBerk.obstacleIdeal, EnvelopeProperty.ProjBerk.mem_obstacleSpace}
\leanok
\uses{def:pn-gauss}
\source{fhx-envelope-20261028:sec:smooth}
Let $d\in\mathbf{Z}_{>0}$ and $\phi\colon(\PP^n)^\an\to\R$. For $m\ge1$ let $V_m(\phi)\subseteq R_{md}[t]$ be the set of
$g$, homogeneous of $x$-degree $md$, with $\log|g|_{\sigma(v)}\le m\,\phi(v)$ for all $v$, and let
$J_m(\phi)\subseteq R[t]$ be the ideal generated by $V_m(\phi)$.
\end{definition}

\begin{lemma}[Properties of $J_m$]\label{lem:bj-Jm-basic}
\lean{EnvelopeProperty.ProjBerk.obstacleIdeal_mul_le, EnvelopeProperty.ProjBerk.obstacleIdeal_eq_span_monomial, EnvelopeProperty.ProjBerk.obstacleIdeal_ne_bot, EnvelopeProperty.ProjBerk.inv_mul_phiIdeal_obstacleIdeal_le, EnvelopeProperty.ProjBerk.isFiniteRatFS_obstacleIdeal}
\leanok
\uses{def:bj-Jm, def:pn-gauss, lem:bj-fs-basic}
Let $\phi$ be bounded. Then: $V_m(\phi)$ is a $k[t]$-submodule; $J_mJ_r\subseteq J_{m+r}$; $J_q^r$ is
generated by bihomogeneous elements of $x$-degree $qrd$; $J_m\ne0$; $J_m$ is bihomogeneous; and
$m^{-1}\phi_{J_m}\in\mathcal H_\Q(d)$ (finite valued) with $m^{-1}\phi_{J_m}\le\phi$.
\end{lemma}

\begin{proof}
The ultrametric inequality, $|t|_{\sigma(v)}=e^{-1}$ and multiplicativity of $\sigma(v)$ give the first
two claims, the second via $V_mV_r\subseteq V_{m+r}$. If $g\in V_m$, its $t$-homogeneous parts lie in $V_m$
by the Gauss formula, so $J_m$ is bihomogeneous and $J_q^r$ is generated by products of $r$ bihomogeneous
elements of $V_q$. If $\phi\ge-C$ then $t^{\lceil mC\rceil}x_i^{md}\in V_m$ for all $i$, so $J_m\ne0$ and
$\phi_{J_m}\ge\max_i\log|x_i^{md}|-\lceil mC\rceil=-\lceil mC\rceil$. By Noetherianity $J_m$ is
generated by finitely many $g_\ell=\sum_\lambda s_{\ell\lambda}t^\lambda\in V_m$; elements of $R[t]$ have
$|\cdot|_{\sigma(v)}\le1$, so $\phi_{J_m}=\max_\ell\log|g_\ell|_{\sigma(\cdot)}
=\max_{\ell,\lambda}(\log|s_{\ell\lambda}|-\lambda)$, a Fubini--Study function for $\mathcal O(md)$ with integer
weights, and $\le m\phi$.
\end{proof}

\begin{proposition}[Envelopes of PL obstacles]\label{prop:bj-envelope-Jm}
\lean{EnvelopeProperty.ProjBerk.envelope_eq_iSup_obstacleIdeal}
\leanok
\uses{lem:bj-fs-envelope, lem:bj-Jm-basic}
\source{fhx-envelope-20261028:prop:ample-envelopes, bj22-trivval:lem:envline}
Let $d\in\mathbf{Z}_{>0}$ and $\phi\colon(\PP^n)^\an\to\R$ bounded, lsc, with $\phi\le0$. Then
$\mathrm P_d(\phi)=\sup_{m\ge1}m^{-1}\phi_{J_m(\phi)}$ pointwise, and $m^{-1}\phi_{J_m}\le
(mr)^{-1}\phi_{J_{mr}}$ for $r\ge1$.
\end{proposition}

\begin{proof}
The monotonicity follows from $J_m^r\subseteq J_{mr}$ and $\phi_{I^r}=r\phi_I$. The functions
$m^{-1}\phi_{J_m}$ are Fubini--Study functions $\le\phi$, hence $\le\mathrm P_d(\phi)$. Conversely, by
Lemma~\ref{lem:bj-fs-envelope} it suffices to show $\psi\le\sup_mm^{-1}\phi_{J_m}$ for
$\psi\in\mathcal H_\Q(d)$, $\psi\le\phi$. Write $\psi=M^{-1}\max_\ell(\log|s_\ell|+\lambda_\ell)$ with
$0\ne s_\ell\in R_{Md}$ and, replacing $M$ by a multiple and $s_\ell$ by powers, $\lambda_\ell\in\mathbf{Z}$.
Evaluating at $\vtriv$ gives $\lambda_\ell\le M\psi(\vtriv)\le M\phi(\vtriv)\le0$. Hence
$g_\ell=s_\ell t^{-\lambda_\ell}\in R_{Md}[t]$ and $\log|g_\ell|_{\sigma(v)}=\log|s_\ell|_v+\lambda_\ell\le
M\psi(v)\le M\phi(v)$, so $g_\ell\in V_M(\phi)$ and $\psi\le M^{-1}\phi_{J_M}$.
\end{proof}

\begin{remark}\label{rmk:bj-test-configurations}
\uses{prop:bj-envelope-Jm}
\cite[Proposition~2.5(2)]{FHX} expresses the envelope of $\phi_{\mathcal L}$, for a test configuration
$(\mathcal X,\mathcal L)$, through the base ideals $\mathfrak a_m$ of $m\mathcal L$ (\cite[Lemma~5.22]{BJ22},
\cite[Lemmas~1.6, 1.7]{BJadd}), and \cite[\S3]{FHX} then passes to the ideals $J_m$ on
$S=X\times\A^1$ with $J_m\cdot\mathcal O_{\mathcal Y}=\mathfrak a_m\mathcal O_{\mathcal Y}(mD)$. By
\cite[Lemma~1.7]{BJadd}, $J_m$ is the ideal of Definition~\ref{def:bj-Jm} for $\phi=\phi_D$. On $\PP^n$ the
proof of Theorem~3.1 only uses the properties in Lemma~\ref{lem:bj-Jm-basic} and
Proposition~\ref{prop:bj-envelope-Jm}, for every PL obstacle, so test configurations
(\cite[Theorem~2.31]{BJ22}) are not needed.
\end{remark}

\subsection{Charts and divisorial points}

\begin{lemma}[Graded extension]\label{lem:bj-graded-extension}
\lean{EnvelopeProperty.ProjBerk.IsGradedNorm, EnvelopeProperty.ProjBerk.ofGradedNorm, EnvelopeProperty.ProjBerk.ofGradedNorm_val_of_isHomogeneous}
\leanok
\uses{def:pn-an}
Let $N$ be a function on the homogeneous elements of $R$ with values in $\R_{\ge0}$ such that $N(0)=0$,
$N(c)=1$ for $c\in k^\times$, $N(fg)=N(f)N(g)$, $N(f+g)\le\max(N(f),N(g))$ for $f,g$ of equal degree, and
$\max_iN(x_i)=1$. Then $|f|\coloneqq\max_dN(f_d)$ is a point of $(\PP^n)^\an$.
\end{lemma}

\begin{proof}
The axioms other than multiplicativity are immediate from those of $N$. Since
$(fg)_c=\sum_{d+e=c}f_dg_e$, we have $N((fg)_c)\le|f||g|$ for all $c$, which settles the case $|f||g|=0$. Otherwise
let $d_0$ (resp.\ $e_0$) be the largest degree with $N(f_{d_0})=|f|$ (resp.\ $N(g_{e_0})=|g|$). In
$(fg)_{d_0+e_0}=f_{d_0}g_{e_0}+\sum f_dg_e$ the other terms have $d>d_0$ or $e>e_0$, hence norm $<|f||g|$, so
$N((fg)_{d_0+e_0})=|f||g|$.
\end{proof}

\begin{lemma}[Affine charts]\label{lem:bj-chart}
\lean{EnvelopeProperty.ProjBerk.chartAlg, EnvelopeProperty.ProjBerk.finiteType_chartAlg, EnvelopeProperty.ProjBerk.chartHomeomorph, EnvelopeProperty.ProjBerk.stdChart, EnvelopeProperty.ProjBerk.mulChart, EnvelopeProperty.ProjBerk.awayComap_mulChart, EnvelopeProperty.ProjBerk.isIntegrallyClosed_chartAlg, EnvelopeProperty.ProjBerk.isIntegrallyClosed_chartAlg_and}
\leanok
\uses{lem:bj-graded-extension, def:berk, lem:principal-open}
\source{fhx-envelope-20261028:sec:normal-local-regularization}
Let $0\ne\sigma\in R_e$, $e\ge1$, and $A_\sigma=(R_\sigma)_0$, the degree-zero part of the localization.
\begin{enumerate}
\item $A_\sigma$ is a finitely generated integrally closed domain, generated by the $x^\alpha/\sigma$ with
$|\alpha|=e$, and every element is $\tau/\sigma^N$ with $\tau\in R_{Ne}$. For $\sigma=x_j$,
$A_{x_j}=k[y^{(j)}]$; for $\sigma=x_0h$ with $h$ the homogenization of $h_1\in k[x_1/x_0,\dots,x_n/x_0]$,
$A_\sigma=k[x_1/x_0,\dots,x_n/x_0]_{h_1}$.
\item $U_\sigma=\{v:|\sigma|_v\ne0\}$ is open, and $\iota_\sigma(v)(\tau/\sigma^N)=|\tau|_v/|\sigma|_v^N$
defines a homeomorphism $\iota_\sigma$ of $U_\sigma$ onto $(\Spec A_\sigma)^\an$.
\item For $f=\tau/\sigma^N\in A_\sigma$, $\log|f|\circ\iota_\sigma=\log|\tau|-N\log|\sigma|$ on $U_\sigma$, and
$\max_{|\alpha|=e}\log|x^\alpha/\sigma|\circ\iota_\sigma=-\log|\sigma|$.
\item $\vtriv\in U_\sigma$ and $\iota_\sigma(\vtriv)=\vtriv$.
\end{enumerate}
\end{lemma}

\begin{proof}
(1) A degree-zero fraction is $\tau/\sigma^N$ with $\deg\tau=Ne$, and $\tau$ is a combination of monomials
$x^\beta$, $|\beta|=Ne$, each a product of $N$ monomials of degree $e$. If $z\in\operatorname{Frac}A_\sigma$
is integral over $A_\sigma$, it is integral over the factorial ring $R_\sigma$, hence lies in
$R_\sigma$; being a quotient of forms of equal degree, it lies in the degree-zero part $A_\sigma$. The two special cases are dehomogenization. (2) $\iota_\sigma(v)$ is a
multiplicative seminorm on $A_\sigma$, trivial on $k$, and depends continuously on $v$. Inverse: for
$y\in(\Spec A_\sigma)^\an$ and $g\in R_c$ put
$N(g)=|g^e/\sigma^c|_y^{1/e}/\mu^c$ with $\mu=\max_{|\alpha|=e}|x^\alpha/\sigma|_y^{1/e}$; here $\mu\ge1$ because
$1=\sum_\alpha c_\alpha x^\alpha/\sigma$, and $|x^\alpha/\sigma|\le\max_i|x_i^e/\sigma|$ because
$(x^\alpha/\sigma)^e=\prod_i(x_i^e/\sigma)^{\alpha_i}$, so $\max_iN(x_i)=1$. Lemma~\ref{lem:bj-graded-extension}
gives a point $v$ with $|\sigma|_v=1/\mu^e\ne0$, and the two constructions are inverse and continuous (the
second as composites of continuous operations on the values $|f|_y$). (3) is the definition of $\iota_\sigma$
together with $\max_{|\alpha|=e}|x^\alpha|_v=1$. (4) $|\sigma|_{\vtriv}=1$.
\end{proof}

\begin{definition}[Divisorial points]\label{def:pn-div}
\lean{EnvelopeProperty.BirationalDVR, EnvelopeProperty.ProjBerk.IsDivisorial}
\leanok
\uses{def:pn-gauss}
\source{fhx-envelope-20261028:subsec:prelim-berkovich, bj22-trivval:sec:Gauss}
A point $v\in(\PP^n)^\an$ is \emph{divisorial} if there are $c\in\Q_{>0}$ and a discrete valuation ring
$R'$ with $k[x_0,\dots,x_n,t]\subseteq R'\subseteq k(x_0,\dots,x_n,t)$, essentially of finite type over $k$,
such that $\log|g|_{\sigma(v)}=-c\operatorname{ord}_{R'}(g)$ for all $g\in R[t]$. We write $(\PP^n)^{\divv}$ for the set of
divisorial points.
\end{definition}

\begin{lemma}[Gauss extension of a discrete valuation ring]\label{lem:bj-gauss-dvr}
\lean{EnvelopeProperty.GaussDVR, EnvelopeProperty.addVal_gaussDVR, EnvelopeProperty.essFiniteType_gaussDVR}
\leanok
Let $V$ be a discrete valuation ring with fraction field $F$, essentially of finite type over $k$, and $X$
an indeterminate. Then $V'=V[X]_{\mathfrak m_VV[X]}$ is a discrete valuation ring with fraction field
$F(X)$, essentially of finite type over $k$, and $\operatorname{ord}_{V'}(\sum_ia_iX^i)=\min_i\operatorname{ord}_V(a_i)$.
\end{lemma}

\begin{proof}
By Gauss's lemma $\nu(\sum a_iX^i)=\min_i\operatorname{ord}_V(a_i)$ is a valuation of $F(X)$ with value group $\mathbf{Z}$; its
valuation ring contains $V[X]$, and its center on $V[X]$ is the prime $\mathfrak m_VV[X]$. A quotient
$p/q$ with $p,q\in V[X]$, $\nu(p)\ge\nu(q)$, can be written with $\nu(q)=0$ after dividing both by a power of
a uniformizer, i.e.\ $q\notin\mathfrak m_VV[X]$; so the valuation ring is $V'$.
\end{proof}

\begin{lemma}[Examples of divisorial points]\label{lem:bj-div-points}
\lean{EnvelopeProperty.ProjBerk.isDivisorial_triv, EnvelopeProperty.ProjBerk.chartDivPoint, EnvelopeProperty.ProjBerk.isDivisorial_chartDivPoint, EnvelopeProperty.ProjBerk.logAbs_gaussExt_chartDivPoint}
\leanok
\uses{def:pn-div, lem:bj-gauss-dvr, lem:bj-graded-extension}
\begin{enumerate}
\item $\vtriv\in(\PP^n)^{\divv}$.
\item Let $0\le j\le n$ and let $V$ be a discrete valuation ring with $B_j\subseteq V\subseteq k(y^{(j)},t)$,
essentially of finite type over $k$, with $b\coloneqq\operatorname{ord}_V(t)>0$ and
$\operatorname{ord}_V(\sum_\lambda f_\lambda t^\lambda)=\min_\lambda(\operatorname{ord}_V(f_\lambda)+\lambda b)$ for $f_\lambda\in k[y^{(j)}]$.
Then $|s|_{v_V}=\exp(-\operatorname{ord}_V(s^{(j)})/b)$ for $s\in R_c$ defines, by graded extension, a divisorial point
$v_V$ with $|x_j|_{v_V}=1$ and $\log|g|_{\sigma(v_V)}=-\operatorname{ord}_V(g^{(j)})/b$ for $g\in R[t]$ homogeneous in $x$.
\end{enumerate}
\end{lemma}

\begin{proof}
(1) $\log|\sum_\lambda s_\lambda t^\lambda|_{\sigma(\vtriv)}=-\min\{\lambda:s_\lambda\ne0\}=-\operatorname{ord}_t$, with
$R'=k[x,t]_{(t)}$ and $c=1$. (2) The hypotheses of Lemma~\ref{lem:bj-graded-extension} hold since
$\operatorname{ord}_V\ge0$ on $k[y^{(j)}]$ and $(x_j)^{(j)}=1$; the Gauss formula for $\sigma(v_V)$ matches the assumed
form of $\operatorname{ord}_V$. Let $R'=V[x_j]_{\mathfrak m_VV[x_j]}\subseteq k(x,t)$ (Lemma~\ref{lem:bj-gauss-dvr},
$x_j$ being transcendental over $k(y^{(j)},t)$). It contains $x_i=y^{(j)}_ix_j$ and $t$. For $g\in R[t]$ with
$x$-homogeneous parts $g_c$, $g=\sum_cg_c^{(j)}x_j^c$, so
$\operatorname{ord}_{R'}(g)=\min_c\operatorname{ord}_V(g_c^{(j)})=-b\max_c\log|g_c|_{\sigma(v_V)}=-b\log|g|_{\sigma(v_V)}$; take $c=1/b$.
\end{proof}

\begin{remark}\label{rmk:bj-div-comparison-BJ}
\uses{def:pn-div}
By \cite[\S1.3]{BJ22}, Gauss extension identifies $X^{\divv}$ in the sense of \cite[\S1.1]{BJ22} with the
$\mathbf G_m$-invariant divisorial valuations $w$ of $X\times\A^1$ with $w(t)=1$. Definition~\ref{def:pn-div}
is the same set (extend $w$ to the affine cone by the Gauss extension in a coordinate $x_j$ of value zero,
as in Lemma~\ref{lem:bj-div-points}). This identification is not used: the uses of $X^{\divv}$ in
\cite[\S3]{FHX} for $X=\PP^n$ are the Frobenius trace estimate at $\sigma(v)$, which needs exactly the
ring $R'$, and the comparison theorem, proved below for Definition~\ref{def:pn-div}.
\end{remark}

\subsection{Rees valuations}

\begin{theorem}[Finiteness of normalization]\label{thm:bj-normalization-finite}
\lean{EnvelopeProperty.finite_of_isIntegralClosure_of_finiteType, EnvelopeProperty.finite_integralClosure}
\leanok
Let $D$ be a finitely generated $k$-algebra which is a domain. The integral closure $\overline D$ of $D$
in its fraction field is a finite $D$-module; in particular it is a finitely generated $k$-algebra and a
Noetherian integrally closed domain.
\end{theorem}

\begin{proof}
By Noether normalization there is a polynomial ring $A=k[z_1,\dots,z_d]\subseteq D$ over which $D$ is finite;
$\overline D$ is the integral closure of $A$ in $L=\operatorname{Frac}D$, a finite extension of
$K=\operatorname{Frac}A$. If $L/K$ is separable, $\overline D$ is finite over the integrally closed Noetherian
ring $A$ (trace form and dual basis). Let $\operatorname{char}k=p>0$ and $L=K(\alpha_1,\dots,\alpha_r)$. The
minimal polynomial of $\alpha_\ell$ is $g_\ell(T^{q_\ell})$ with $g_\ell\in K[T]$ separable and $q_\ell$ a power
of $p$; let $q=\max_\ell q_\ell$. In an algebraic closure $\Omega$ of $L$ let $K_q=k(z_1^{1/q},\dots,z_d^{1/q})$
and $A_q=k[z_1^{1/q},\dots,z_d^{1/q}]$, a polynomial ring finite free over $A$ (monomial basis). Since $k$ is
perfect, the coefficients of $g_\ell$ have $q_\ell$-th roots in $K_q$, and $g_\ell(T^{q_\ell})=h_\ell(T)^{q_\ell}$
with $h_\ell\in K_q[T]$ the image of $g_\ell$ under the field isomorphism $K\to K^{1/q_\ell}$,
$c\mapsto c^{1/q_\ell}$; $h_\ell$ is separable and $h_\ell(\alpha_\ell)=0$. Hence $L'=K_q(\alpha_1,\dots,\alpha_r)$ is
finite separable over $K_q$, and the integral closure $C'$ of $A_q$ in $L'$ is finite over $A_q$, hence over
$A$. Every element of $L$ integral over $A$ is integral over $A_q$, so $\overline D\subseteq C'$ is an
$A$-submodule of a finite module over the Noetherian ring $A$, hence finite.
(Formalized route: with $K_s$ the separable closure of $K$ in $L$, $L/K_s$ is purely inseparable of some exponent $n$ and $x\mapsto x^{p^n}$ maps $\overline D$ injectively into the integral closure of $A$ in $K_s$, compatibly with the Frobenius of $A$, which is finite since $k$ is perfect.)
\end{proof}

\begin{lemma}[Associated primes of principal ideals in normal rings]\label{lem:bj-dvr-minimal}
\lean{EnvelopeProperty.isDiscreteValuationRing_localization_of_mem_minimalPrimes, EnvelopeProperty.associatedPrimes_eq_minimalPrimes, EnvelopeProperty.mem_span_singleton_iff_forall_minimalPrimes_dvd}
\leanok
Let $E$ be a Noetherian integrally closed domain and $0\ne h\in E$. For every associated prime $P$ of
$E/hE$, $E_P$ is a discrete valuation ring and $P$ is a minimal prime of $hE$. Consequently the
associated primes of $E/hE$ are the finitely many minimal primes of $hE$, the same for all powers of $h$,
and for $f\in E$: $f\in hE$ if and only if $\operatorname{ord}_{E_P}(f)\ge\operatorname{ord}_{E_P}(h)$ for every minimal prime $P$ of $hE$.
\end{lemma}

\begin{proof}
Let $P=(hE:a)$ with $a\notin hE$, and $z=a/h$, so $zP\subseteq E$ and $zPE_P\subseteq E_P$. If
$zPE_P\subseteq PE_P$, the determinant trick applied to the finite faithful $E_P$-module $PE_P$ shows that $z$
is integral over $E_P$, which is integrally closed; so $a\in hE_P$, i.e.\ $sa\in hE$ for some $s\notin P$,
contradicting $(hE:a)=P$. Hence $zPE_P=E_P$, $PE_P=z^{-1}E_P$ is principal and nonzero, and the
Noetherian local domain $E_P$ is a discrete valuation ring; so $\dim E_P=1$ and $P$ is minimal over $hE$.
Minimal primes are associated, and $hE$, $h^rE$ have the same minimal primes. Finally, if $f\notin hE$,
the nonzero module $(fE+hE)/hE$ has an associated prime $P=(hE:bf)$, which is associated to $E/hE$; if
$f\in hE_P$ then $sbf\in hE$ for some $s\notin P$, a contradiction. So $f\notin hE_P$ for that $P$.
\end{proof}

\begin{definition}[Rees valuation rings]\label{def:bj-rees}
\lean{EnvelopeProperty.reesChart, EnvelopeProperty.reesNormalChart, EnvelopeProperty.ReesIndex, EnvelopeProperty.reesValuationRing, EnvelopeProperty.isDiscreteValuationRing_reesValuationRing, EnvelopeProperty.finite_reesIndex, EnvelopeProperty.essFiniteType_reesValuationRing}
\leanok
\uses{thm:bj-normalization-finite}
\source{bj22-trivval:defi:Rees}
Let $B$ be a finitely generated $k$-domain and $I=(g_1,\dots,g_s)$ with all $g_\ell\ne0$. Let
$D_\ell=B[g_1/g_\ell,\dots,g_s/g_\ell]\subseteq\operatorname{Frac}B$ and $E_\ell$ its integral closure. The
\emph{Rees valuation rings} $\mathrm{RV}(I;g)$ are the localizations $(E_\ell)_P$, $1\le\ell\le s$, $P$ a
minimal prime of $g_\ell E_\ell$. By Theorem~\ref{thm:bj-normalization-finite} and
Lemma~\ref{lem:bj-dvr-minimal} this is a finite set of discrete valuation rings containing $B$, essentially
of finite type over $k$, and $\operatorname{ord}_V(I)=\operatorname{ord}_V(g_\ell)$ for $V=(E_\ell)_P$ (as $g_i/g_\ell\in V$).
\end{definition}

\begin{theorem}[Rees valuations detect integral closure]\label{thm:bj-rees}
\lean{EnvelopeProperty.valuation_le_pow_of_forall_dvd, EnvelopeProperty.valuation_le_sup_pow_of_forall_dvd, EnvelopeProperty.isGreatest_valuation_span_pow, EnvelopeProperty.exists_dvd_iff_addVal_le}
\leanok
\uses{def:bj-rees, lem:bj-dvr-minimal, thm:bj-normalization-finite}
\source{bj22-trivval:lem:CLT, bj22-trivval:lem:Shilov}
In the setting of Definition~\ref{def:bj-rees}, let $w$ be a multiplicative ultrametric seminorm on $B$ with
$|b|_w\le1$ for $b\in B$ and $\max_\ell|g_\ell|_w>0$. Let $F\in B$ and $r\in\N$ with
$\operatorname{ord}_V(F)\ge r\operatorname{ord}_V(I)$ for all $V\in\mathrm{RV}(I;g)$. Then $|F|_w\le\max_\ell|g_\ell|_w^r
=\sup_{g\in I^r}|g|_w$.
\end{theorem}

\begin{proof}
The last equality holds because $|b|_w\le1$ on $B$. Choose $\ell$ with $|g_\ell|_w$ maximal, hence $>0$.
Then $w$ extends to $B_{g_\ell}$ by $|b/g_\ell^m|=|b|_w/|g_\ell|_w^m$, and is $\le1$ on $D_\ell$, which is
generated over $B$ by the $g_i/g_\ell$ of norm $\le1$. By hypothesis and Definition~\ref{def:bj-rees},
$\operatorname{ord}_V(F)\ge r\operatorname{ord}_V(g_\ell)$ for every minimal prime of $g_\ell E_\ell$, which are those of $g_\ell^rE_\ell$;
Lemma~\ref{lem:bj-dvr-minimal} gives $z\coloneqq F/g_\ell^r\in E_\ell$. Thus $z$ satisfies
$z^N+c_1z^{N-1}+\dots+c_N=0$ with $c_i\in D_\ell$. If $|z|>1$, then $|c_iz^{N-i}|\le|z|^{N-i}<|z|^N$ for
all $i$, contradicting the ultrametric inequality; so $|F|_w\le|g_\ell|_w^r$.
\end{proof}

\begin{lemma}[Invariance under the torus]\label{lem:bj-rees-invariant}
\lean{EnvelopeProperty.addVal_X_pos_reesValuationRing, EnvelopeProperty.addVal_reesValuationRing_eq_inf}
\leanok
\uses{def:bj-rees,algebra:rootable-action-fixes-minimal-primes}
Let $B=k[y_1,\dots,y_n][t]$ and $I=(g_1,\dots,g_s)$ with $g_\ell=h_\ell t^{e_\ell}$, $0\ne h_\ell\in k[y]$, and
$t^N\in I$ for some $N$. For every $V\in\mathrm{RV}(I;g)$: $\operatorname{ord}_V(t)>0$, and
$\operatorname{ord}_V(\sum_\lambda f_\lambda t^\lambda)=\min_\lambda(\operatorname{ord}_V(f_\lambda)+\lambda\operatorname{ord}_V(t))$ for $f_\lambda\in k[y]$.
\end{lemma}

\begin{proof}
Let $V=(E_\ell)_P$. Since $t^N\in I$, $t^N\in g_\ell D_\ell\subseteq P$, so $t\in P$. For $\lambda\in k^\times$
the $k[y]$-automorphism $\alpha_\lambda\colon t\mapsto\lambda t$ maps $g_i/g_\ell$ to $\lambda^{e_i-e_\ell}g_i/g_\ell$,
hence preserves $D_\ell$, $E_\ell$ and $g_\ell E_\ell$, and permutes the finite set of minimal primes of
$g_\ell E_\ell$. The stabilizer of $P$ has finite index in $k^\times$, which is divisible since $k$ is
algebraically closed, so it is all of $k^\times$; thus $\operatorname{ord}_V\circ\alpha_\lambda=\operatorname{ord}_V$. For
$f=\sum_{\lambda\le L}f_\lambda t^\lambda$ choose distinct $\mu_0,\dots,\mu_L\in k^\times$; inverting the
Vandermonde matrix expresses each $f_\lambda t^\lambda$ as a $k$-combination of the $\alpha_{\mu_i}(f)$, so
$\operatorname{ord}_V(f_\lambda t^\lambda)\ge\operatorname{ord}_V(f)$. The reverse inequality is the ultrametric inequality.
\end{proof}

\subsection{The Shilov lemma and divisorial comparison}

\begin{lemma}[Finite Fubini--Study functions]\label{lem:bj-finite-fs}
\lean{EnvelopeProperty.ProjBerk.X_mem_radical_span, EnvelopeProperty.ProjBerk.X_pow_mem_span_dehomGen, EnvelopeProperty.ProjBerk.IsFiniteRatFS.exists_coneSup}
\leanok
\uses{lem:bj-graded-extension, def:pn-gauss}
\source{bj22-trivval:prop:FSideal}
Let $\psi=m^{-1}\max_i(\log|s_i|+\lambda_i)\in\mathcal H_\Q(b)$ with $0\ne s_i\in R_{mb}$, $\lambda_i\in\mathbf{Z}_{\le0}$,
$\max_i\lambda_i=0$. Then $\mathfrak m^e\subseteq(s_i)_i$ for some $e$; the ideal $I\subseteq R[t]$ generated by
the $g_i=s_it^{-\lambda_i}$ satisfies $\phi_I=m\psi$, $x_j^et^{N}\in I$ for $N=\max_i(-\lambda_i)$ and all $j$, and
hence $t^N\in I^{(j)}$; $I^{(j)}$ is generated by the $g_i^{(j)}$.
\end{lemma}

\begin{proof}
If the $s_i$ had a common zero, i.e.\ a homogeneous prime $\mathfrak p\supseteq(s_i)$ with $\mathfrak p\not\supseteq\mathfrak m$
(Nullstellensatz), the trivial semivaluation of $\mathfrak p$ ($|f|=0$ for $f\in\mathfrak p$ homogeneous, $1$
otherwise, extended by Lemma~\ref{lem:bj-graded-extension}) would be a point where $\psi=-\infty$. Hence
$\sqrt{(s_i)}\supseteq\mathfrak m$. Then $\phi_I=\max_i\log|g_i|_{\sigma(\cdot)}=m\psi$, and
$x^\alpha=\sum_ic_is_i$ gives $x^\alpha t^N=\sum_ic_it^{N+\lambda_i}g_i$ for $|\alpha|=e$.
\end{proof}

\begin{lemma}[Shilov boundary]\label{lem:bj-shilov}
\lean{EnvelopeProperty.ProjBerk.IsFiniteRatFS.exists_finset_iSup_sub_eq, EnvelopeProperty.ProjBerk.isDivisorial_reesPoint}
\leanok
\uses{lem:bj-finite-fs, thm:bj-rees, lem:bj-rees-invariant, lem:bj-div-points}
\source{bj22-trivval:lem:Shilov}
For every $\psi\in\mathcal H_\Q(b)$ there is a finite set $\Sigma\subseteq(\PP^n)^{\divv}$ containing $\vtriv$
such that for every $a\in\Q_{\ge0}$ and every $\varphi\in\mathcal H^{\mathrm{gf}}_\Q(a)$,
$\sup_{(\PP^n)^\an}(\varphi-\psi)=\max_\Sigma(\varphi-\psi)$.
\end{lemma}

\begin{proof}
Subtracting a rational constant from $\psi$, and replacing $m$ by a multiple and the $s_i$ by powers, we may
assume $\psi$ is as in Lemma~\ref{lem:bj-finite-fs};
let $I$, $g_i$, $N$ be as there. For each $j$, Lemma~\ref{lem:bj-rees-invariant} and
Lemma~\ref{lem:bj-div-points} attach to every $V\in\mathrm{RV}(I^{(j)};g^{(j)})$ a divisorial point $v_V$;
let $\Sigma$ consist of $\vtriv$ and these points. Write $\varphi=M^{-1}\max_\ell(\log|\sigma_\ell|+\mu_\ell)$ with
$0\ne\sigma_\ell\in R_{Ma}$, $\mu_\ell\in\mathbf{Z}$, and let $c=\max_\Sigma(\varphi-\psi)$. Since $\varphi(\vtriv)$ is
finite and $\psi(\vtriv)=0$, $c\ge\varphi(\vtriv)\ge\mu_\ell/M$ for all $\ell$; the values at points of $\Sigma$
are rational or $-\infty$, so $c\in\Q$. Choose $r\ge1$ with $rMc\in\mathbf{Z}$ and put
$a_\ell=m(rMc-r\mu_\ell)\in\N$, $F_\ell=\sigma_\ell^{rm}t^{a_\ell}$. Multiplying by $rMm$ shows that, at any $v$,
$(\varphi-\psi)(v)\le c$ is equivalent to $\log|F_\ell|_{\sigma(v)}\le rM\phi_I(v)=\phi_{I^{rM}}(v)$ for all $\ell$.
At $v=v_V\in\Sigma$ (chart $j$), $\log|F_\ell|_{\sigma(v_V)}=-\operatorname{ord}_V(F_\ell^{(j)})/\operatorname{ord}_V(t)$ and
$\phi_{I^{rM}}(v_V)=-rM\operatorname{ord}_V(I^{(j)})/\operatorname{ord}_V(t)$; so the hypothesis $(\varphi-\psi)(v_V)\le c$ says
$\operatorname{ord}_V(F_\ell^{(j)})\ge rM\operatorname{ord}_V(I^{(j)})$ for every $V\in\mathrm{RV}(I^{(j)};g^{(j)})$. Now let
$v\in(\PP^n)^\an$ and choose $j$ with $|x_j|_v=1$. Then $w(f)=|x_j^cf|_{\sigma(v)}$ ($c$ at least the degree of $f$ in $y^{(j)}$; independent of $c$ as $|x_j|_v=1$) is a
multiplicative seminorm on $B_j$, $\le1$ on $B_j$, with $|t|_w=e^{-1}$, hence
$\max_i|g_i^{(j)}|_w\ge|t^N|_w>0$. Theorem~\ref{thm:bj-rees} gives
$\log|F_\ell|_{\sigma(v)}=\log|F_\ell^{(j)}|_w\le\phi_{I^{rM}}(v)$.
\end{proof}

\begin{lemma}[Suprema against PL functions]\label{lem:bj-sup-pl}
\lean{EnvelopeProperty.ProjBerk.IsRatPL.exists_finset_iSup_sub_eq}
\leanok
\uses{lem:bj-shilov, thm:bj-regularization}
\source{bj22-trivval:lem:supPL}
For every $\psi\in\operatorname{PL}$ there is a finite $\Sigma\subseteq(\PP^n)^{\divv}$ such that
$\sup(u-\psi)=\max_\Sigma(u-\psi)$ for every $\theta\ge0$ and every $u\in\PSH(\theta)$.
\end{lemma}

\begin{proof}
Write $\psi=\psi_1-\psi_2$ with $\psi_\nu\in\mathcal H_\Q(b_\nu)$ and take $\Sigma$ for $\psi_1$ from
Lemma~\ref{lem:bj-shilov}. Let $u_\alpha\in\mathcal H_\Q(a_\alpha)$ decrease to $u$
(Theorem~\ref{thm:bj-regularization}). Then $u_\alpha+\psi_2\in\mathcal H^{\mathrm{gf}}_\Q(a_\alpha+b_2)$, so for every $v$,
$(u-\psi)(v)\le(u_\alpha+\psi_2-\psi_1)(v)\le\max_\Sigma(u_\alpha-\psi)$. As $\Sigma$ is finite and $u_\alpha$
decreases to $u$ pointwise, the right side tends to $\max_\Sigma(u-\psi)$.
\end{proof}

\begin{theorem}[Divisorial comparison]\label{thm:bj-divisorial-comparison}
\lean{EnvelopeProperty.ProjBerk.IsPSH.le_of_forall_isDivisorial_le, EnvelopeProperty.ProjBerk.IsPSH.eq_of_forall_isDivisorial_eq}
\leanok
\uses{lem:bj-sup-pl, thm:bj-pl-dense, lem:bj-usc-pl}
\source{bj22-trivval:thm:suppsh, bj22-trivval:cor:Hausdorff, fhx-envelope-20261028:prop:BJ-package}
Let $\theta\ge0$, $u\in\PSH(\theta)$ and $h\colon(\PP^n)^\an\to[-\infty,\infty)$ usc. If $u\le h$ on
$(\PP^n)^{\divv}$, then $u\le h$ everywhere. In particular a $\theta$-psh function is determined by its
restriction to $(\PP^n)^{\divv}$.
\end{theorem}

\begin{proof}
Let $\psi\in\operatorname{PL}$ with $\psi>h$. By Lemma~\ref{lem:bj-sup-pl}, $\sup(u-\psi)=\max_\Sigma(u-\psi)\le0$
since $\Sigma\subseteq(\PP^n)^{\divv}$. So $u\le\psi$ for all such $\psi$, and $u\le h$ by
Lemma~\ref{lem:bj-usc-pl}. For the last claim apply this to two psh functions in both directions.
\end{proof}

\begin{proposition}[Increasing nets]\label{prop:bj-ep-criterion}
\lean{EnvelopeProperty.ProjBerk.hasEnvelopeProperty_of_forall_monotone, EnvelopeProperty.ProjBerk.IsPSH.eq_uscReg_iSup}
\leanok
\uses{thm:bj-divisorial-comparison}
\source{fhx-envelope-20261028:prop:EP-criterion, bj22-trivval:thm:envppty}
Let $\theta\ge0$. If every bounded-above increasing net $(u_i)$ in $\PSH(\theta)$ admits $u\in\PSH(\theta)$ with
$u=\sup_iu_i$ on $(\PP^n)^{\divv}$, then $\theta$ has the envelope property. For any such net and $u$,
$u=(\sup_iu_i)^\star$.
\end{proposition}

\begin{proof}
Second claim: $u_i\le u$ on $(\PP^n)^{\divv}$, hence everywhere (Theorem~\ref{thm:bj-divisorial-comparison}),
so $(\sup u_i)^\star\le u$; and $u=\sup u_i\le(\sup u_i)^\star$ on $(\PP^n)^{\divv}$, hence everywhere. For the
envelope property, given a bounded-above family $(u_i)_{i\in I}$, apply the hypothesis to the increasing net
$\max_{i\in J}u_i$ over finite $J\subseteq I$ (finite maxima are psh); its limit $u$ is
$(\sup_J\max_{i\in J}u_i)^\star=(\sup_iu_i)^\star$.
\end{proof}

\begin{remark}\label{rmk:bj-ep-criterion-converse}
\uses{prop:bj-ep-criterion}
The converse implication of \cite[Proposition~2.4]{FHX}, envelope property $\Rightarrow$ convergence of
increasing nets on $X^{\divv}$, needs that divisorial points are nonnegligible \cite[Theorem~5.6]{BJ22},
i.e.\ $(\sup_iu_i)^\star=\sup_iu_i$ on $X^{\divv}$. Its proof (\cite[Lemma~5.7]{BJ22}) realizes a divisorial
point as a Rees valuation of a flag ideal (\cite[Lemma~2.15]{BJ22}, via \cite[Lemma~1.20]{BJ22}), uses PL
interpolation at Rees valuations (\cite[Lemmas~2.18, 2.19]{BJ22}, via relative ampleness on the normalized
blowup) and the finiteness of psh functions at divisorial points (\cite[Proposition~1.28,
Corollary~4.17]{BJ22}, via intersection numbers on a normal model). This direction is used in
\cite[\S\S4, 6--8]{FHX} for general $X$ (\S\S4, 6, 7 in revision \texttt{832e086}), but neither in \cite[\S3]{FHX} nor in \cite[\S5]{FHX} for $X=\PP^n$.
The finiteness of psh functions on $X^{\divv}$ in \cite[Proposition~2.1(2)]{FHX} is likewise not used on
$\PP^n$: Section~5 only uses finiteness at $\vtriv$.
\end{remark}

\subsection{Theorem 3.1 for projective space}

\begin{remark}[Assembly of Theorem~3.1 for $\PP^n$ in characteristic $p$]\label{rmk:bj-pn-ep-assembly}
\uses{lem:bj-envppty, prop:bj-envelope-criterion, lem:bj-env-basic, prop:bj-envelope-Jm, lem:bj-Jm-basic, lem:test-ideals, lem:pn-divisorial-trace, def:pn-div, lem:bj-sup-pl, thm:bj-psh2, thm:bj-psh-basic}
Theorem~\ref{thm:pn-ep} follows from the nodes above together with Lemmas~\ref{lem:test-ideals} and
\ref{lem:pn-divisorial-trace}, as in \cite[\S3]{FHX}. (a) By Lemma~\ref{lem:bj-envppty} with rationals
$\theta_j\downarrow\theta$, $\theta_j>\theta$, it suffices to treat $\theta=a\in\Q_{>0}$; by
Proposition~\ref{prop:bj-envelope-criterion} it suffices that $\mathrm P_a(\phi)$ is psh for $\phi\in\operatorname{PL}$;
by Lemma~\ref{lem:bj-env-basic}, replacing $(a,\phi)$ by $(ma,m\phi)$ and subtracting a rational constant
$\ge\max\phi$, we may assume $a=d\in\mathbf{Z}_{>0}$ and $\phi\le0$, $\phi\in\operatorname{PL}$. (b) Put $J_m=J_m(\phi)$ and
$\mathrm Q=\mathrm P_d(\phi)=\sup_mm^{-1}\phi_{J_m}$ (Proposition~\ref{prop:bj-envelope-Jm}). By
Lemma~\ref{lem:bj-Jm-basic}, $(J_m)$ is a graded sequence of nonzero bihomogeneous ideals with $J_q$ generated in
$x$-degree $\le qd$, so Lemma~\ref{lem:test-ideals} gives asymptotic test ideals $B_m=\tau(J_q^{m/q})$ ($q$
divisible) with $J_m\subseteq B_m$, $B_{m\ell}\subseteq B_m^\ell$, $B_m$ generated in $x$-degree $\le md$.
Padding generators with monomials (whose $\max_i$ of $\log|x_i^c|$ is $0$), $u_m=m^{-1}\phi_{B_m}\in
\mathcal H^{\mathrm{gf}}_\Q(d)$. (c) $q^{-1}\phi_{J_q}\le(qm)^{-1}\phi_{J_{qm}}\le u_{qm}\le u_m$, so $\mathrm Q\le u_m\le0$;
$u_{(j+1)!}\le u_{j!}$, and $u=\lim_ju_{j!}\in\PSH(d)$ by PSH2 since $u\ge\mathrm Q\ge\phi_{J_1}>-\infty$. (d) Let
$v\in(\PP^n)^{\divv}$ with $R'$, $c$ as in Definition~\ref{def:pn-div}. Lemma~\ref{lem:pn-divisorial-trace}
(trace estimate for the Frobenius trace of $k[x,t]$ at $\operatorname{ord}_{R'}$, with a constant $A$ depending on $R'$)
and $B_m=\Phi_e(J_q^{\lceil mp^e/q\rceil})$, $e\gg0$, give $\operatorname{ord}_{R'}(B_m)\ge(m/q)\operatorname{ord}_{R'}(J_q)-A$, i.e.\
$u_m(v)\le q^{-1}\phi_{J_q}(v)+cA/m\le\mathrm Q(v)+cA/m$; hence $u=\mathrm Q$ on $(\PP^n)^{\divv}$. (e) Thus $u\le\phi$
on $(\PP^n)^{\divv}$; Lemma~\ref{lem:bj-sup-pl} with $\psi=\phi$ gives $u\le\phi$ everywhere, so
$u\le\mathrm P_d(\phi)=\mathrm Q\le u$, and $\mathrm P_d(\phi)=u$ is psh. In (d) the test ideals of
Lemma~\ref{lem:test-ideals} live in the polynomial ring $k[x_0,\dots,x_n,t]$, the Frobenius trace being that of
the variables, which form a $p$-basis of $k(x,t)$; this is why Definition~\ref{def:pn-div} uses a discrete
valuation ring of $k(x,t)$ containing $k[x,t]$.
\end{remark}

\begin{remark}[Characteristic zero]\label{rmk:bj-char-zero}
\uses{rmk:bj-pn-ep-assembly}
For $\operatorname{char}k=0$, \cite[\S3]{FHX} cites \cite[Theorem~5.20]{BJ22}, whose proof resolves the test
configuration and uses multiplier ideals with their subadditivity and uniform global generation. On $\PP^n$
the argument of Remark~\ref{rmk:bj-pn-ep-assembly} applies verbatim with the asymptotic multiplier ideals of
$(J_m)$ on the smooth variety $\PP^n\times\A^1$ in place of test ideals, provided one has: $J_m\subseteq\mathcal
J(J_\bullet^m)$, subadditivity, generation of $\mathcal J(J_\bullet^m)$ in $x$-degree $\le md$ (since
$K_{\PP^n}+(n+1)H=0$), and the valuative bound $w(\mathcal J(I^c))\ge c\,w(I)-A(w)$ at divisorial $w$. Everything
in this chapter except Lemma~\ref{lem:pn-divisorial-trace} and Lemma~\ref{lem:test-ideals} is
characteristic free. The route planned for the formalization in characteristic $0$ (Section on characteristic zero,
Theorem~\ref{thm:c0-pn-ep}) avoids resolution and multiplier ideals: an $L^2$ (Ohsawa--Takegoshi, H\"ormander)
bound replaces subadditivity and global generation; the choice is pending the author's decision.
\end{remark}

\subsection{The compactification lemma on projective space}

In this subsection $P=\PP^n$, $0\ne\sigma\in R_e$ with $e\ge1$, $S=P_\sigma$, $A_\sigma=k[S]$ and
$\iota_\sigma$ are as in Lemma~\ref{lem:bj-chart}, and we identify $S^\an$ with $U_\sigma$ through
$\iota_\sigma$. For $c\in\Q_{>0}$ the class $c\,\mathcal O_P(e)$ is $ce\in\Q$.

\begin{lemma}[Compactification, part (1)]\label{lem:bj-compactification-regularize}
\lean{EnvelopeProperty.ProjBerk.IsPSH.exists_chartFS, EnvelopeProperty.ProjBerk.IsPSH.isRegularizable_chart, EnvelopeProperty.ProjBerk.IsPSH.isRegularizable_mulChart}
\leanok
\uses{lem:bj-chart, cor:bj-countable-regularization, def:fs, lem:bj-fs-basic}
\source{fhx-envelope-20261028:lem:compactification}
Assume $k$ countable. Let $c\in\Q_{>0}$ and $\Phi\in\PSH(ce)$. There are $u_j\in\FS(S)$ decreasing to
$\Phi-c\log|\sigma|$ on $S^\an$ with $u_j+c\log|\sigma|$ bounded. In particular
$(\Phi-c\log|\sigma|)|_{S^\an}\in\Rcal(S)$.
\end{lemma}

\begin{proof}
By Corollary~\ref{cor:bj-countable-regularization} (fixed class) take $\Phi_j\in\mathcal H_\Q(ce)$ decreasing to
$\Phi$. Write $\Phi_j=M^{-1}\max_\ell(\log|\tau_\ell|+\lambda_\ell)$ with $\tau_\ell\in R_{Mce}$ and, after replacing $M$
by $eM$ and $\tau_\ell$ by $\tau_\ell^e$, $Mc\in\N$. Then $f_\ell=\tau_\ell/\sigma^{Mc}\in A_\sigma$ and
$\Phi_j-c\log|\sigma|=M^{-1}\max_\ell(\log|f_\ell|+\lambda_\ell)$ on $S^\an$ (Lemma~\ref{lem:bj-chart}(3)). Choose a rational
$C_j\ge-\min\Phi_j$, so that $\Phi_j-c\log|\sigma|\ge-c\log|\sigma|-C_j$. With the generators $z_\alpha=x^\alpha/\sigma$
of $A_\sigma$, $\log|z_\alpha|\le-\log|\sigma|$, so the terms $c\log|z_\alpha|-C_j$ and the constant $-C_j$ are
$\le\Phi_j-c\log|\sigma|$ ($\log|\sigma|\le0$). Hence
$u_j=\max\{M^{-1}\max_\ell(\log|f_\ell|+\lambda_\ell),\max_\alpha(c\log|z_\alpha|-C_j),-C_j\}\in\FS(S)$ equals
$\Phi_j-c\log|\sigma|$; the $u_j$ decrease, and $u_j+c\log|\sigma|=\Phi_j$ is bounded. The limit is not
identically $-\infty$ because $\vtriv\in S^\an$ and $\Phi(\vtriv)>-\infty$.
\end{proof}

\begin{lemma}[Divisibility from growth]\label{lem:bj-divisibility}
\lean{EnvelopeProperty.ProjBerk.pow_floor_dvd_of_logAbs_le, EnvelopeProperty.ProjBerk.pow_floor_dvd_of_log_le}
\leanok
\uses{lem:bj-graded-extension}
Let $\tau\in R_{Ne}$, $\beta\ge0$ and $C\in\R$ with $\log|\tau|-N\log|\sigma|\le-\beta\log|\sigma|+C$ on $U_\sigma$.
If $N\ge\beta$, then $\sigma^{\lfloor N-\beta\rfloor}$ divides $\tau$ in $R$.
\end{lemma}

\begin{proof}
Factor $\sigma=\gamma\prod_\iota\sigma_\iota^{m_\iota}$ into homogeneous primes of the factorial ring $R$. For $s>0$,
$|f|_{s,\iota}=\exp(-s\operatorname{ord}_{\sigma_\iota}(f))$ on homogeneous $f$ satisfies the hypotheses of
Lemma~\ref{lem:bj-graded-extension} (some $x_i$ is not divisible by $\sigma_\iota$, as $n\ge1$), and the point lies
in $U_\sigma$. The hypothesis there reads $-s\operatorname{ord}_{\sigma_\iota}(\tau)+sNm_\iota\le s\beta m_\iota+C$; dividing by $s$ and
letting $s\to\infty$ gives $\operatorname{ord}_{\sigma_\iota}(\tau)\ge(N-\beta)m_\iota\ge\lfloor N-\beta\rfloor m_\iota$ for all $\iota$.
\end{proof}

\begin{lemma}[Compactification, part (2), regularizable functions]\label{lem:bj-compactification-extend}
\lean{EnvelopeProperty.ProjBerk.isPSH_chartExt, EnvelopeProperty.ProjBerk.isPSH_chartExt_of_fs, EnvelopeProperty.ProjBerk.isPSH_mulChartExt, EnvelopeProperty.ProjBerk.isPSH_chartExt'}
\leanok
\uses{lem:bj-chart, lem:bj-divisibility, thm:import-growth-control, lem:reg-triv-finite, def:fs, lem:bj-fs-basic}
\source{fhx-envelope-20261028:lem:compactification}
Let $w\in\Rcal(S)$, $\kappa\in\Q_{>0}$ and $B_0\in\R$ with $w\le-\kappa\log|\sigma|+B_0$ on $S^\an$. Assume the imported
Theorem~\ref{thm:import-growth-control}. For every rational $a>\kappa$, the function $\phi$ equal to
$w+a\log|\sigma|$ on $S^\an$ and to $-\infty$ elsewhere lies in $\PSH(ae)$.
\end{lemma}

\begin{proof}
Choose rationals $\kappa'\in[\kappa,a)$ and $c\in(1,a/\kappa')$. The function
\[
 g=\max\{\max_{|\alpha|=e}\kappa'\log|x^\alpha/\sigma|,0\}\in\FS(S)
\]
equals $-\kappa'\log|\sigma|$
(Lemma~\ref{lem:bj-chart}(3)), and $w\le g+B_0$. Theorem~\ref{thm:import-growth-control} gives $u_j\in\FS(S)$
decreasing to $w$ with $u_j\le-\alpha\log|\sigma|+C_j$, $\alpha=c\kappa'<a$. Write
$u_j=\max\{\max_i(a_i\log|f_i|+b_i),b_0\}$ with $a_i\in\Q_{>0}$ and $f_i=\tau_i/\sigma^{N_i}$, $\tau_i\in R_{N_ie}$
(Lemma~\ref{lem:bj-chart}(1)). Each term satisfies $a_i\log|f_i|\le-\alpha\log|\sigma|+C'_j$. Fix a positive integer
$r$ with $\alpha+a_i/r<a$ for all $i$. If $N_i\ge\alpha/a_i$, Lemma~\ref{lem:bj-divisibility} applied to $\tau_i^r$
(with $N=rN_i$, $\beta=r\alpha/a_i$) gives $\tau_i^r=\sigma^{N'_i}\tau'_i$ with $N'_i=\lfloor r(N_i-\alpha/a_i)\rfloor$, so
$a_i\log|f_i|=(a_i/r)(\log|\tau'_i|-(rN_i-N'_i)\log|\sigma|)$ with $rN_i-N'_i\le r\alpha/a_i+1$, hence
$(a_i/r)(rN_i-N'_i)\le\alpha+a_i/r<a$; if $N_i<\alpha/a_i$, then $a_iN_i<\alpha<a$ and we keep $\tau'_i=\tau_i$ with
$r$ replaced by $1$. In all cases
$a_i\log|f_i|+b_i+a\log|\sigma|=\epsilon_i\log|\tau'_i|+\delta_i\log|\sigma|+b_i$ with $\epsilon_i\in\Q_{>0}$, $\delta_i\in\Q_{>0}$ and
$\epsilon_i\deg\tau'_i+\delta_ie=ae$. Taking a common denominator $M$, the function
\[
 \Psi_j=\max\{\max_i(\epsilon_i\log|\tau'_i|+\delta_i\log|\sigma|+b_i),a\log|\sigma|+b_0\}
\]
is a Fubini--Study function for
$\mathcal O(ae)$; it equals $u_j+a\log|\sigma|$ on $U_\sigma$ and $-\infty$ on its complement, since every term has a
positive power of $\sigma$. The $\Psi_j$ decrease to $\phi$, and $\phi(\vtriv)=w(\vtriv)>-\infty$
(Lemma~\ref{lem:reg-triv-finite}); hence $\phi\in\PSH(ae)$ by Theorem~\ref{thm:bj-psh2}.
\end{proof}

\begin{remark}\label{rmk:bj-compactification-full}
\uses{lem:bj-compactification-regularize, lem:bj-compactification-extend}
\cite[Lemma~5.10(2)]{FHX} allows any $u$ that is, near every compact set, a decreasing limit of convex PL functions;
its proof glues truncations and needs \cite[Lemma~5.5]{FHX} (locally $L$-convex functions are $L$-psh), whose proof
uses \cite[Theorem~2.31, Definition~3.2, Lemma~4.3]{BJ22} (PL functions as nef test configurations, with relative
ampleness and Serre vanishing on the test configuration) and the positivity criterion of Chambert-Loir--Ducros
\cite{CLD12}. Every use of \cite[Lemma~5.10(2)]{FHX} outside the proof of \cite[Lemma~5.11(5),(6)]{FHX} has a
regularizable input: the extension step of Proposition~\ref{prop:nonflat-norm} (on $\PP^n$, with
$\sigma=x_0\cdot h$ and $S\subseteq\A^n$ via Lemma~\ref{lem:bj-chart}(1)), and \cite[Corollary~5.26]{FHX} (on a general
normal projective variety). For the former, Lemma~\ref{lem:bj-compactification-extend} suffices, conditionally on the
already imported Theorem~\ref{thm:import-growth-control}; the latter needs the same argument on a general normal
projective variety, hence the global theory of \cite{BJ22} beyond $\PP^n$ (Lemma~\ref{lem:pg-compactification-extend}). Discharging that import, or Lemma~5.10(2) in full, requires
\cite[Lemma~5.5]{FHX}.
\end{remark}
